\section{研究背景与问题提出}

\subsection{研究背景}


\subsection{现有研究进展与不足}



\subsection{问题提出与研究切入点}

首先需要在正常运行条件下识别西安轨道交通系统中的韧性薄弱站点、区间和换乘节点，明确网络中哪些位置在常态下已经承受较大压力或具有较高脆弱性。其次，需要在天气扰动情境下刻画站点客流、时段分布与 OD 联系强度的变化，识别外部冲击对轨道出行需求的作用方式。最后，还需要进一步分析这些客流变化如何经由网络联系传导到站点、区间与换乘节点的运行压力，并考察站点环境特征在这一过程中发挥的加剧或缓冲作用。

在这一研究框架中，站点周边 POI 和建成环境特征并不作为独立研究对象，而是作为解释变量和调节变量嵌入分析过程，用于揭示不同类型站点对天气扰动的响应差异。也就是说，天气是外部冲击源，客流变化是关键中介过程，系统韧性变化是研究结果，而站点环境特征则帮助解释这一传导链为何在不同空间单元上表现出不同强度和不同方向。

\subsection{研究意义}

本研究有助于推动轨道交通韧性研究由静态结构识别向动态扰动传导分析拓展，将天气扰动、客流响应和网络韧性纳入同一分析框架之中，并强化站点环境异质性在韧性研究中的解释作用。从实践层面看，本研究能够为西安轨道交通系统识别重点关注站点、区间与换乘节点提供依据，也有助于为恶劣天气条件下的客流组织、运力调配、换乘疏导和精细化管理提供决策参考。尤其是在城市轨道网络持续扩展、站点功能日趋多样的背景下，围绕天气扰动下韧性变化机制开展研究，具有较强的现实针对性和应用价值。
